% =====================================================================
%  Společná preambule textových dokumentů 4EK453 (A4): příklady ke cvičení
%  (cviceni/) a otázky ke zkoušce (zkouska/). Vkládá se za \documentclass:
%    \documentclass[11pt,a4paper]{article}
%    % =====================================================================
%  Společná preambule textových dokumentů 4EK453 (A4): příklady ke cvičení
%  (cviceni/) a otázky ke zkoušce (zkouska/). Vkládá se za \documentclass:
%    \documentclass[11pt,a4paper]{article}
%    % =====================================================================
%  Společná preambule textových dokumentů 4EK453 (A4): příklady ke cvičení
%  (cviceni/) a otázky ke zkoušce (zkouska/). Vkládá se za \documentclass:
%    \documentclass[11pt,a4paper]{article}
%    % =====================================================================
%  Společná preambule textových dokumentů 4EK453 (A4): příklady ke cvičení
%  (cviceni/) a otázky ke zkoušce (zkouska/). Vkládá se za \documentclass:
%    \documentclass[11pt,a4paper]{article}
%    \input{../prezentace/src/_shared/dokument.tex}
%  Překlad: xelatex (2×). Značení a makra jsou stejná jako v prezentacích
%  (makra.tex), barvy stejné jako v beamerthemeFIS.sty.
% =====================================================================

% Pořadí: babel a mathtools před unicode-math (jinak kolize aktivního znaku „-“
% a mathtools přepíše \underbrace verzí pro staré fonty)
\usepackage[czech]{babel}
\usepackage{amsmath}
\usepackage{mathtools}
% Písma: Latin Modern, volaná jménem souboru (jako v _quarto.yml)
\usepackage{fontspec}
\usepackage{unicode-math}
\setmainfont{lmroman10}[Extension=.otf, UprightFont=*-regular, BoldFont=*-bold,
  ItalicFont=*-italic, BoldItalicFont=*-bolditalic]
\setsansfont{lmsans10}[Extension=.otf, UprightFont=*-regular, BoldFont=*-bold,
  ItalicFont=*-oblique, BoldItalicFont=*-boldoblique]
\setmonofont{lmmono10}[Extension=.otf, UprightFont=*-regular, ItalicFont=*-italic]
\setmathfont{latinmodern-math.otf}

\AtBeginDocument{\shorthandoff{-"}}

\usepackage[a4paper, left=25mm, right=25mm, top=25mm, bottom=25mm, headsep=8mm]{geometry}
\usepackage{xcolor}
\usepackage{booktabs}
\usepackage{array}
\usepackage{enumitem}
\usepackage{comment}
\usepackage{fancyvrb}
\usepackage{titlesec}
\usepackage{fancyhdr}
\usepackage{etoolbox}
\usepackage[hidelinks]{hyperref}

% Barvy FIS VŠE (stejné jako v beamerthemeFIS.sty)
\definecolor{fisZelena}{HTML}{00957D}
\definecolor{fisZelenaTmava}{HTML}{006E5C}
\definecolor{fisZelenaSvetla}{HTML}{DFF1EC}
\definecolor{fisSeda}{HTML}{54585A}
\definecolor{fisSedaSvetla}{HTML}{7C8183}
\definecolor{fisLinka}{HTML}{DCE0DF}
\definecolor{fisPodklad}{HTML}{EBF0EE}
\definecolor{fisCerna}{HTML}{1D1F20}
\definecolor{fisModra}{HTML}{0072B0}
\definecolor{fisCervena}{HTML}{D0453A}

% cesta ke sdíleným souborům; dokument hlouběji ve stromu si ji nastaví před \input,
% např. \newcommand{\sdilene}{../../prezentace/src/_shared/}
\providecommand{\sdilene}{../prezentace/src/_shared/}
% název kurzu v~titulku a záhlaví; dokument pro více kurzů (MP) si před \input
% nastaví \def\kurz{} a řádek s~kurzem se nevysází
\providecommand{\kurz}{4EK453 Principy ekonometrie}
\input{\sdilene makra.tex}

% ---------------------------------------------------------------------
%  Sazba
% ---------------------------------------------------------------------
\setlength{\parindent}{0pt}
\setlength{\parskip}{4pt plus 1pt}
\renewcommand{\arraystretch}{1.1}
\frenchspacing
\setlength{\emergencystretch}{2em}% dlouhé \texttt{…} a vzorce raději řídčeji než přes okraj

\titleformat{\section}{\sffamily\Large\bfseries\color{fisZelenaTmava}}{\thesection}{0.6em}{}
\titleformat{\subsection}{\sffamily\large\bfseries\color{fisSeda}}{\thesubsection}{0.6em}{}
\titlespacing*{\section}{0pt}{16pt plus 4pt}{6pt}
\titlespacing*{\subsection}{0pt}{12pt plus 3pt}{4pt}
\renewcommand{\paragraph}[1]{\par\medskip\noindent{\sffamily\bfseries\color{fisZelenaTmava}#1}\quad}

% Titulek: \title, \subtitle, \author, \institution; datum = datum překladu (verze)
\makeatletter
\newcommand{\subtitle}[1]{\gdef\@subtitle{#1}}
\newcommand{\institution}[1]{\gdef\@institution{#1}}
\gdef\@subtitle{}\gdef\@institution{}
\renewcommand{\maketitle}{%
  \thispagestyle{plain}%
  {\color{fisZelena}\rule{\linewidth}{2pt}}\par\vspace{4pt}%
  \ifx\kurz\@empty\else{\sffamily\small\color{fisSeda}\kurz\par}\vspace{2pt}\fi
  {\sffamily\LARGE\bfseries\color{fisCerna}\@title\par}\vspace{2pt}%
  \ifx\@subtitle\@empty\else{\sffamily\large\color{fisSeda}\@subtitle\par}\fi
  \vspace{6pt}%
  {\small\@author\ \textbar\ \@institution\hfill verze z~\@date\par}%
  \vspace{2pt}{\color{fisZelena}\rule{\linewidth}{0.6pt}}\par\vspace{8pt}}
\makeatother
\date{\today}

\pagestyle{fancy}
\fancyhf{}
\renewcommand{\headrulewidth}{0.4pt}
\renewcommand{\headrule}{{\color{fisLinka}\hrule height\headrulewidth}}
\makeatletter
\fancyhead[L]{\sffamily\footnotesize\color{fisSeda}\kurz}
\fancyhead[R]{\sffamily\footnotesize\color{fisSeda}\@title\ \textbar\ \@date}
\makeatother
\fancyfoot[C]{\sffamily\footnotesize\color{fisSeda}\thepage}
\fancypagestyle{plain}{\fancyhf{}\renewcommand{\headrulewidth}{0pt}%
  \fancyfoot[C]{\sffamily\footnotesize\color{fisSeda}\thepage}}

% ---------------------------------------------------------------------
%  Příklady
% ---------------------------------------------------------------------
% \begin{priklad} … \end{priklad}         úloha na papír
% \begin{priklad}[R] … \end{priklad}      úloha na počítači (v R)
% \koumaci                                označení „(Pro koumáky.)“
% \begin{prikladlist} \item … \end{prikladlist}   podúlohy a), b), …
% \begin{reseni} … \end{reseni}           řešení; tisk řídí \includecomment / \excludecomment
\newcounter{priklad}[section]
\renewcommand{\thepriklad}{\ifnum\value{section}>0 \thesection.\fi\arabic{priklad}}
% značka R: malý obdélník bez \strut, spodkem barvy na účaří (nezvětšuje výšku řádku)
\newcommand{\znackaR}{\,{\setlength{\fboxsep}{1.3pt}\raisebox{\fboxsep}{\colorbox{fisZelena}{\sffamily\scriptsize\bfseries\color{white}\,R\,}}}}
% (Pro koumáky.) = vyžaduje aparát, který studenti nemusí mít (integrály, Choleského
% rozklad, obtížnější důkazy); u podúloh je dáváme na konec příkladu
\newcommand{\koumaci}{(\textit{Pro koumáky.})\ }% závorky stojatě
\NewDocumentEnvironment{priklad}{o}{%
  \par\addvspace{8pt}\refstepcounter{priklad}%
  \noindent{\sffamily\bfseries\color{fisZelenaTmava}Příklad~\thepriklad}%
  \IfValueT{#1}{\hspace{0.3em}\znackaR}\quad\ignorespaces}{\par\addvspace{4pt}}
\newlist{prikladlist}{enumerate}{1}
\setlist[prikladlist]{label=\textsl{\alph*})~, leftmargin=2em, itemsep=1pt, topsep=2pt}
\newcommand{\rada}{\par\smallskip\textit{Rada:} }
\newcommand{\pozn}{\par\smallskip\textit{Poznámka:} }

% Číslovaný seznam otázek s předponou (T1, T2, … / P1, P2, …), číslování
% pokračuje přes nadpisy: \begin{otazky}{T} … \end{otazky}
\newlist{otazkyT}{enumerate}{1}
\setlist[otazkyT]{label=\textbf{\sffamily\color{fisZelenaTmava}T\arabic*.}, ref=T\arabic*, resume,
  leftmargin=3em, itemsep=3pt, topsep=3pt}
\newlist{otazkyP}{enumerate}{1}
\setlist[otazkyP]{label=\textbf{\sffamily\color{fisZelenaTmava}P\arabic*.}, ref=P\arabic*, resume,
  leftmargin=3em, itemsep=6pt, topsep=3pt}

% Výpisy z R (vstupují ze souborů vygenerovaných skriptem, nikdy ne opsané ručně)
\fvset{fontsize=\small, xleftmargin=0.5em, frame=leftline, rulecolor=\color{fisLinka},
  framerule=1.5pt, framesep=5pt}
% výpis se nedělí mezi stránky
\newcommand{\vystupR}[1]{\par\noindent\begin{minipage}{\linewidth}\VerbatimInput{#1}\end{minipage}\par}

%  Překlad: xelatex (2×). Značení a makra jsou stejná jako v prezentacích
%  (makra.tex), barvy stejné jako v beamerthemeFIS.sty.
% =====================================================================

% Pořadí: babel a mathtools před unicode-math (jinak kolize aktivního znaku „-“
% a mathtools přepíše \underbrace verzí pro staré fonty)
\usepackage[czech]{babel}
\usepackage{amsmath}
\usepackage{mathtools}
% Písma: Latin Modern, volaná jménem souboru (jako v _quarto.yml)
\usepackage{fontspec}
\usepackage{unicode-math}
\setmainfont{lmroman10}[Extension=.otf, UprightFont=*-regular, BoldFont=*-bold,
  ItalicFont=*-italic, BoldItalicFont=*-bolditalic]
\setsansfont{lmsans10}[Extension=.otf, UprightFont=*-regular, BoldFont=*-bold,
  ItalicFont=*-oblique, BoldItalicFont=*-boldoblique]
\setmonofont{lmmono10}[Extension=.otf, UprightFont=*-regular, ItalicFont=*-italic]
\setmathfont{latinmodern-math.otf}

\AtBeginDocument{\shorthandoff{-"}}

\usepackage[a4paper, left=25mm, right=25mm, top=25mm, bottom=25mm, headsep=8mm]{geometry}
\usepackage{xcolor}
\usepackage{booktabs}
\usepackage{array}
\usepackage{enumitem}
\usepackage{comment}
\usepackage{fancyvrb}
\usepackage{titlesec}
\usepackage{fancyhdr}
\usepackage{etoolbox}
\usepackage[hidelinks]{hyperref}

% Barvy FIS VŠE (stejné jako v beamerthemeFIS.sty)
\definecolor{fisZelena}{HTML}{00957D}
\definecolor{fisZelenaTmava}{HTML}{006E5C}
\definecolor{fisZelenaSvetla}{HTML}{DFF1EC}
\definecolor{fisSeda}{HTML}{54585A}
\definecolor{fisSedaSvetla}{HTML}{7C8183}
\definecolor{fisLinka}{HTML}{DCE0DF}
\definecolor{fisPodklad}{HTML}{EBF0EE}
\definecolor{fisCerna}{HTML}{1D1F20}
\definecolor{fisModra}{HTML}{0072B0}
\definecolor{fisCervena}{HTML}{D0453A}

% cesta ke sdíleným souborům; dokument hlouběji ve stromu si ji nastaví před \input,
% např. \newcommand{\sdilene}{../../prezentace/src/_shared/}
\providecommand{\sdilene}{../prezentace/src/_shared/}
% název kurzu v~titulku a záhlaví; dokument pro více kurzů (MP) si před \input
% nastaví \def\kurz{} a řádek s~kurzem se nevysází
\providecommand{\kurz}{4EK453 Principy ekonometrie}
\input{\sdilene makra.tex}

% ---------------------------------------------------------------------
%  Sazba
% ---------------------------------------------------------------------
\setlength{\parindent}{0pt}
\setlength{\parskip}{4pt plus 1pt}
\renewcommand{\arraystretch}{1.1}
\frenchspacing
\setlength{\emergencystretch}{2em}% dlouhé \texttt{…} a vzorce raději řídčeji než přes okraj

\titleformat{\section}{\sffamily\Large\bfseries\color{fisZelenaTmava}}{\thesection}{0.6em}{}
\titleformat{\subsection}{\sffamily\large\bfseries\color{fisSeda}}{\thesubsection}{0.6em}{}
\titlespacing*{\section}{0pt}{16pt plus 4pt}{6pt}
\titlespacing*{\subsection}{0pt}{12pt plus 3pt}{4pt}
\renewcommand{\paragraph}[1]{\par\medskip\noindent{\sffamily\bfseries\color{fisZelenaTmava}#1}\quad}

% Titulek: \title, \subtitle, \author, \institution; datum = datum překladu (verze)
\makeatletter
\newcommand{\subtitle}[1]{\gdef\@subtitle{#1}}
\newcommand{\institution}[1]{\gdef\@institution{#1}}
\gdef\@subtitle{}\gdef\@institution{}
\renewcommand{\maketitle}{%
  \thispagestyle{plain}%
  {\color{fisZelena}\rule{\linewidth}{2pt}}\par\vspace{4pt}%
  \ifx\kurz\@empty\else{\sffamily\small\color{fisSeda}\kurz\par}\vspace{2pt}\fi
  {\sffamily\LARGE\bfseries\color{fisCerna}\@title\par}\vspace{2pt}%
  \ifx\@subtitle\@empty\else{\sffamily\large\color{fisSeda}\@subtitle\par}\fi
  \vspace{6pt}%
  {\small\@author\ \textbar\ \@institution\hfill verze z~\@date\par}%
  \vspace{2pt}{\color{fisZelena}\rule{\linewidth}{0.6pt}}\par\vspace{8pt}}
\makeatother
\date{\today}

\pagestyle{fancy}
\fancyhf{}
\renewcommand{\headrulewidth}{0.4pt}
\renewcommand{\headrule}{{\color{fisLinka}\hrule height\headrulewidth}}
\makeatletter
\fancyhead[L]{\sffamily\footnotesize\color{fisSeda}\kurz}
\fancyhead[R]{\sffamily\footnotesize\color{fisSeda}\@title\ \textbar\ \@date}
\makeatother
\fancyfoot[C]{\sffamily\footnotesize\color{fisSeda}\thepage}
\fancypagestyle{plain}{\fancyhf{}\renewcommand{\headrulewidth}{0pt}%
  \fancyfoot[C]{\sffamily\footnotesize\color{fisSeda}\thepage}}

% ---------------------------------------------------------------------
%  Příklady
% ---------------------------------------------------------------------
% \begin{priklad} … \end{priklad}         úloha na papír
% \begin{priklad}[R] … \end{priklad}      úloha na počítači (v R)
% \koumaci                                označení „(Pro koumáky.)“
% \begin{prikladlist} \item … \end{prikladlist}   podúlohy a), b), …
% \begin{reseni} … \end{reseni}           řešení; tisk řídí \includecomment / \excludecomment
\newcounter{priklad}[section]
\renewcommand{\thepriklad}{\ifnum\value{section}>0 \thesection.\fi\arabic{priklad}}
% značka R: malý obdélník bez \strut, spodkem barvy na účaří (nezvětšuje výšku řádku)
\newcommand{\znackaR}{\,{\setlength{\fboxsep}{1.3pt}\raisebox{\fboxsep}{\colorbox{fisZelena}{\sffamily\scriptsize\bfseries\color{white}\,R\,}}}}
% (Pro koumáky.) = vyžaduje aparát, který studenti nemusí mít (integrály, Choleského
% rozklad, obtížnější důkazy); u podúloh je dáváme na konec příkladu
\newcommand{\koumaci}{(\textit{Pro koumáky.})\ }% závorky stojatě
\NewDocumentEnvironment{priklad}{o}{%
  \par\addvspace{8pt}\refstepcounter{priklad}%
  \noindent{\sffamily\bfseries\color{fisZelenaTmava}Příklad~\thepriklad}%
  \IfValueT{#1}{\hspace{0.3em}\znackaR}\quad\ignorespaces}{\par\addvspace{4pt}}
\newlist{prikladlist}{enumerate}{1}
\setlist[prikladlist]{label=\textsl{\alph*})~, leftmargin=2em, itemsep=1pt, topsep=2pt}
\newcommand{\rada}{\par\smallskip\textit{Rada:} }
\newcommand{\pozn}{\par\smallskip\textit{Poznámka:} }

% Číslovaný seznam otázek s předponou (T1, T2, … / P1, P2, …), číslování
% pokračuje přes nadpisy: \begin{otazky}{T} … \end{otazky}
\newlist{otazkyT}{enumerate}{1}
\setlist[otazkyT]{label=\textbf{\sffamily\color{fisZelenaTmava}T\arabic*.}, ref=T\arabic*, resume,
  leftmargin=3em, itemsep=3pt, topsep=3pt}
\newlist{otazkyP}{enumerate}{1}
\setlist[otazkyP]{label=\textbf{\sffamily\color{fisZelenaTmava}P\arabic*.}, ref=P\arabic*, resume,
  leftmargin=3em, itemsep=6pt, topsep=3pt}

% Výpisy z R (vstupují ze souborů vygenerovaných skriptem, nikdy ne opsané ručně)
\fvset{fontsize=\small, xleftmargin=0.5em, frame=leftline, rulecolor=\color{fisLinka},
  framerule=1.5pt, framesep=5pt}
% výpis se nedělí mezi stránky
\newcommand{\vystupR}[1]{\par\noindent\begin{minipage}{\linewidth}\VerbatimInput{#1}\end{minipage}\par}

%  Překlad: xelatex (2×). Značení a makra jsou stejná jako v prezentacích
%  (makra.tex), barvy stejné jako v beamerthemeFIS.sty.
% =====================================================================

% Pořadí: babel a mathtools před unicode-math (jinak kolize aktivního znaku „-“
% a mathtools přepíše \underbrace verzí pro staré fonty)
\usepackage[czech]{babel}
\usepackage{amsmath}
\usepackage{mathtools}
% Písma: Latin Modern, volaná jménem souboru (jako v _quarto.yml)
\usepackage{fontspec}
\usepackage{unicode-math}
\setmainfont{lmroman10}[Extension=.otf, UprightFont=*-regular, BoldFont=*-bold,
  ItalicFont=*-italic, BoldItalicFont=*-bolditalic]
\setsansfont{lmsans10}[Extension=.otf, UprightFont=*-regular, BoldFont=*-bold,
  ItalicFont=*-oblique, BoldItalicFont=*-boldoblique]
\setmonofont{lmmono10}[Extension=.otf, UprightFont=*-regular, ItalicFont=*-italic]
\setmathfont{latinmodern-math.otf}

\AtBeginDocument{\shorthandoff{-"}}

\usepackage[a4paper, left=25mm, right=25mm, top=25mm, bottom=25mm, headsep=8mm]{geometry}
\usepackage{xcolor}
\usepackage{booktabs}
\usepackage{array}
\usepackage{enumitem}
\usepackage{comment}
\usepackage{fancyvrb}
\usepackage{titlesec}
\usepackage{fancyhdr}
\usepackage{etoolbox}
\usepackage[hidelinks]{hyperref}

% Barvy FIS VŠE (stejné jako v beamerthemeFIS.sty)
\definecolor{fisZelena}{HTML}{00957D}
\definecolor{fisZelenaTmava}{HTML}{006E5C}
\definecolor{fisZelenaSvetla}{HTML}{DFF1EC}
\definecolor{fisSeda}{HTML}{54585A}
\definecolor{fisSedaSvetla}{HTML}{7C8183}
\definecolor{fisLinka}{HTML}{DCE0DF}
\definecolor{fisPodklad}{HTML}{EBF0EE}
\definecolor{fisCerna}{HTML}{1D1F20}
\definecolor{fisModra}{HTML}{0072B0}
\definecolor{fisCervena}{HTML}{D0453A}

% cesta ke sdíleným souborům; dokument hlouběji ve stromu si ji nastaví před \input,
% např. \newcommand{\sdilene}{../../prezentace/src/_shared/}
\providecommand{\sdilene}{../prezentace/src/_shared/}
% název kurzu v~titulku a záhlaví; dokument pro více kurzů (MP) si před \input
% nastaví \def\kurz{} a řádek s~kurzem se nevysází
\providecommand{\kurz}{4EK453 Principy ekonometrie}
\input{\sdilene makra.tex}

% ---------------------------------------------------------------------
%  Sazba
% ---------------------------------------------------------------------
\setlength{\parindent}{0pt}
\setlength{\parskip}{4pt plus 1pt}
\renewcommand{\arraystretch}{1.1}
\frenchspacing
\setlength{\emergencystretch}{2em}% dlouhé \texttt{…} a vzorce raději řídčeji než přes okraj

\titleformat{\section}{\sffamily\Large\bfseries\color{fisZelenaTmava}}{\thesection}{0.6em}{}
\titleformat{\subsection}{\sffamily\large\bfseries\color{fisSeda}}{\thesubsection}{0.6em}{}
\titlespacing*{\section}{0pt}{16pt plus 4pt}{6pt}
\titlespacing*{\subsection}{0pt}{12pt plus 3pt}{4pt}
\renewcommand{\paragraph}[1]{\par\medskip\noindent{\sffamily\bfseries\color{fisZelenaTmava}#1}\quad}

% Titulek: \title, \subtitle, \author, \institution; datum = datum překladu (verze)
\makeatletter
\newcommand{\subtitle}[1]{\gdef\@subtitle{#1}}
\newcommand{\institution}[1]{\gdef\@institution{#1}}
\gdef\@subtitle{}\gdef\@institution{}
\renewcommand{\maketitle}{%
  \thispagestyle{plain}%
  {\color{fisZelena}\rule{\linewidth}{2pt}}\par\vspace{4pt}%
  \ifx\kurz\@empty\else{\sffamily\small\color{fisSeda}\kurz\par}\vspace{2pt}\fi
  {\sffamily\LARGE\bfseries\color{fisCerna}\@title\par}\vspace{2pt}%
  \ifx\@subtitle\@empty\else{\sffamily\large\color{fisSeda}\@subtitle\par}\fi
  \vspace{6pt}%
  {\small\@author\ \textbar\ \@institution\hfill verze z~\@date\par}%
  \vspace{2pt}{\color{fisZelena}\rule{\linewidth}{0.6pt}}\par\vspace{8pt}}
\makeatother
\date{\today}

\pagestyle{fancy}
\fancyhf{}
\renewcommand{\headrulewidth}{0.4pt}
\renewcommand{\headrule}{{\color{fisLinka}\hrule height\headrulewidth}}
\makeatletter
\fancyhead[L]{\sffamily\footnotesize\color{fisSeda}\kurz}
\fancyhead[R]{\sffamily\footnotesize\color{fisSeda}\@title\ \textbar\ \@date}
\makeatother
\fancyfoot[C]{\sffamily\footnotesize\color{fisSeda}\thepage}
\fancypagestyle{plain}{\fancyhf{}\renewcommand{\headrulewidth}{0pt}%
  \fancyfoot[C]{\sffamily\footnotesize\color{fisSeda}\thepage}}

% ---------------------------------------------------------------------
%  Příklady
% ---------------------------------------------------------------------
% \begin{priklad} … \end{priklad}         úloha na papír
% \begin{priklad}[R] … \end{priklad}      úloha na počítači (v R)
% \koumaci                                označení „(Pro koumáky.)“
% \begin{prikladlist} \item … \end{prikladlist}   podúlohy a), b), …
% \begin{reseni} … \end{reseni}           řešení; tisk řídí \includecomment / \excludecomment
\newcounter{priklad}[section]
\renewcommand{\thepriklad}{\ifnum\value{section}>0 \thesection.\fi\arabic{priklad}}
% značka R: malý obdélník bez \strut, spodkem barvy na účaří (nezvětšuje výšku řádku)
\newcommand{\znackaR}{\,{\setlength{\fboxsep}{1.3pt}\raisebox{\fboxsep}{\colorbox{fisZelena}{\sffamily\scriptsize\bfseries\color{white}\,R\,}}}}
% (Pro koumáky.) = vyžaduje aparát, který studenti nemusí mít (integrály, Choleského
% rozklad, obtížnější důkazy); u podúloh je dáváme na konec příkladu
\newcommand{\koumaci}{(\textit{Pro koumáky.})\ }% závorky stojatě
\NewDocumentEnvironment{priklad}{o}{%
  \par\addvspace{8pt}\refstepcounter{priklad}%
  \noindent{\sffamily\bfseries\color{fisZelenaTmava}Příklad~\thepriklad}%
  \IfValueT{#1}{\hspace{0.3em}\znackaR}\quad\ignorespaces}{\par\addvspace{4pt}}
\newlist{prikladlist}{enumerate}{1}
\setlist[prikladlist]{label=\textsl{\alph*})~, leftmargin=2em, itemsep=1pt, topsep=2pt}
\newcommand{\rada}{\par\smallskip\textit{Rada:} }
\newcommand{\pozn}{\par\smallskip\textit{Poznámka:} }

% Číslovaný seznam otázek s předponou (T1, T2, … / P1, P2, …), číslování
% pokračuje přes nadpisy: \begin{otazky}{T} … \end{otazky}
\newlist{otazkyT}{enumerate}{1}
\setlist[otazkyT]{label=\textbf{\sffamily\color{fisZelenaTmava}T\arabic*.}, ref=T\arabic*, resume,
  leftmargin=3em, itemsep=3pt, topsep=3pt}
\newlist{otazkyP}{enumerate}{1}
\setlist[otazkyP]{label=\textbf{\sffamily\color{fisZelenaTmava}P\arabic*.}, ref=P\arabic*, resume,
  leftmargin=3em, itemsep=6pt, topsep=3pt}

% Výpisy z R (vstupují ze souborů vygenerovaných skriptem, nikdy ne opsané ručně)
\fvset{fontsize=\small, xleftmargin=0.5em, frame=leftline, rulecolor=\color{fisLinka},
  framerule=1.5pt, framesep=5pt}
% výpis se nedělí mezi stránky
\newcommand{\vystupR}[1]{\par\noindent\begin{minipage}{\linewidth}\VerbatimInput{#1}\end{minipage}\par}

%  Překlad: xelatex (2×). Značení a makra jsou stejná jako v prezentacích
%  (makra.tex), barvy stejné jako v beamerthemeFIS.sty.
% =====================================================================

% Pořadí: babel a mathtools před unicode-math (jinak kolize aktivního znaku „-“
% a mathtools přepíše \underbrace verzí pro staré fonty)
\usepackage[czech]{babel}
\usepackage{amsmath}
\usepackage{mathtools}
% Písma: Latin Modern, volaná jménem souboru (jako v _quarto.yml)
\usepackage{fontspec}
\usepackage{unicode-math}
\setmainfont{lmroman10}[Extension=.otf, UprightFont=*-regular, BoldFont=*-bold,
  ItalicFont=*-italic, BoldItalicFont=*-bolditalic]
\setsansfont{lmsans10}[Extension=.otf, UprightFont=*-regular, BoldFont=*-bold,
  ItalicFont=*-oblique, BoldItalicFont=*-boldoblique]
\setmonofont{lmmono10}[Extension=.otf, UprightFont=*-regular, ItalicFont=*-italic]
\setmathfont{latinmodern-math.otf}

\AtBeginDocument{\shorthandoff{-"}}

\usepackage[a4paper, left=25mm, right=25mm, top=25mm, bottom=25mm, headsep=8mm]{geometry}
\usepackage{xcolor}
\usepackage{booktabs}
\usepackage{array}
\usepackage{enumitem}
\usepackage{comment}
\usepackage{fancyvrb}
\usepackage{titlesec}
\usepackage{fancyhdr}
\usepackage{etoolbox}
\usepackage[hidelinks]{hyperref}

% Barvy FIS VŠE (stejné jako v beamerthemeFIS.sty)
\definecolor{fisZelena}{HTML}{00957D}
\definecolor{fisZelenaTmava}{HTML}{006E5C}
\definecolor{fisZelenaSvetla}{HTML}{DFF1EC}
\definecolor{fisSeda}{HTML}{54585A}
\definecolor{fisSedaSvetla}{HTML}{7C8183}
\definecolor{fisLinka}{HTML}{DCE0DF}
\definecolor{fisPodklad}{HTML}{EBF0EE}
\definecolor{fisCerna}{HTML}{1D1F20}
\definecolor{fisModra}{HTML}{0072B0}
\definecolor{fisCervena}{HTML}{D0453A}

% cesta ke sdíleným souborům; dokument hlouběji ve stromu si ji nastaví před \input,
% např. \newcommand{\sdilene}{../../prezentace/src/_shared/}
\providecommand{\sdilene}{../prezentace/src/_shared/}
% název kurzu v~titulku a záhlaví; dokument pro více kurzů (MP) si před \input
% nastaví \def\kurz{} a řádek s~kurzem se nevysází
\providecommand{\kurz}{4EK453 Principy ekonometrie}
\input{\sdilene makra.tex}

% ---------------------------------------------------------------------
%  Sazba
% ---------------------------------------------------------------------
\setlength{\parindent}{0pt}
\setlength{\parskip}{4pt plus 1pt}
\renewcommand{\arraystretch}{1.1}
\frenchspacing
\setlength{\emergencystretch}{2em}% dlouhé \texttt{…} a vzorce raději řídčeji než přes okraj

\titleformat{\section}{\sffamily\Large\bfseries\color{fisZelenaTmava}}{\thesection}{0.6em}{}
\titleformat{\subsection}{\sffamily\large\bfseries\color{fisSeda}}{\thesubsection}{0.6em}{}
\titlespacing*{\section}{0pt}{16pt plus 4pt}{6pt}
\titlespacing*{\subsection}{0pt}{12pt plus 3pt}{4pt}
\renewcommand{\paragraph}[1]{\par\medskip\noindent{\sffamily\bfseries\color{fisZelenaTmava}#1}\quad}

% Titulek: \title, \subtitle, \author, \institution; datum = datum překladu (verze)
\makeatletter
\newcommand{\subtitle}[1]{\gdef\@subtitle{#1}}
\newcommand{\institution}[1]{\gdef\@institution{#1}}
\gdef\@subtitle{}\gdef\@institution{}
\renewcommand{\maketitle}{%
  \thispagestyle{plain}%
  {\color{fisZelena}\rule{\linewidth}{2pt}}\par\vspace{4pt}%
  \ifx\kurz\@empty\else{\sffamily\small\color{fisSeda}\kurz\par}\vspace{2pt}\fi
  {\sffamily\LARGE\bfseries\color{fisCerna}\@title\par}\vspace{2pt}%
  \ifx\@subtitle\@empty\else{\sffamily\large\color{fisSeda}\@subtitle\par}\fi
  \vspace{6pt}%
  {\small\@author\ \textbar\ \@institution\hfill verze z~\@date\par}%
  \vspace{2pt}{\color{fisZelena}\rule{\linewidth}{0.6pt}}\par\vspace{8pt}}
\makeatother
\date{\today}

\pagestyle{fancy}
\fancyhf{}
\renewcommand{\headrulewidth}{0.4pt}
\renewcommand{\headrule}{{\color{fisLinka}\hrule height\headrulewidth}}
\makeatletter
\fancyhead[L]{\sffamily\footnotesize\color{fisSeda}\kurz}
\fancyhead[R]{\sffamily\footnotesize\color{fisSeda}\@title\ \textbar\ \@date}
\makeatother
\fancyfoot[C]{\sffamily\footnotesize\color{fisSeda}\thepage}
\fancypagestyle{plain}{\fancyhf{}\renewcommand{\headrulewidth}{0pt}%
  \fancyfoot[C]{\sffamily\footnotesize\color{fisSeda}\thepage}}

% ---------------------------------------------------------------------
%  Příklady
% ---------------------------------------------------------------------
% \begin{priklad} … \end{priklad}         úloha na papír
% \begin{priklad}[R] … \end{priklad}      úloha na počítači (v R)
% \koumaci                                označení „(Pro koumáky.)“
% \begin{prikladlist} \item … \end{prikladlist}   podúlohy a), b), …
% \begin{reseni} … \end{reseni}           řešení; tisk řídí \includecomment / \excludecomment
\newcounter{priklad}[section]
\renewcommand{\thepriklad}{\ifnum\value{section}>0 \thesection.\fi\arabic{priklad}}
% značka R: malý obdélník bez \strut, spodkem barvy na účaří (nezvětšuje výšku řádku)
\newcommand{\znackaR}{\,{\setlength{\fboxsep}{1.3pt}\raisebox{\fboxsep}{\colorbox{fisZelena}{\sffamily\scriptsize\bfseries\color{white}\,R\,}}}}
% (Pro koumáky.) = vyžaduje aparát, který studenti nemusí mít (integrály, Choleského
% rozklad, obtížnější důkazy); u podúloh je dáváme na konec příkladu
\newcommand{\koumaci}{(\textit{Pro koumáky.})\ }% závorky stojatě
\NewDocumentEnvironment{priklad}{o}{%
  \par\addvspace{8pt}\refstepcounter{priklad}%
  \noindent{\sffamily\bfseries\color{fisZelenaTmava}Příklad~\thepriklad}%
  \IfValueT{#1}{\hspace{0.3em}\znackaR}\quad\ignorespaces}{\par\addvspace{4pt}}
\newlist{prikladlist}{enumerate}{1}
\setlist[prikladlist]{label=\textsl{\alph*})~, leftmargin=2em, itemsep=1pt, topsep=2pt}
\newcommand{\rada}{\par\smallskip\textit{Rada:} }
\newcommand{\pozn}{\par\smallskip\textit{Poznámka:} }

% Číslovaný seznam otázek s předponou (T1, T2, … / P1, P2, …), číslování
% pokračuje přes nadpisy: \begin{otazky}{T} … \end{otazky}
\newlist{otazkyT}{enumerate}{1}
\setlist[otazkyT]{label=\textbf{\sffamily\color{fisZelenaTmava}T\arabic*.}, ref=T\arabic*, resume,
  leftmargin=3em, itemsep=3pt, topsep=3pt}
\newlist{otazkyP}{enumerate}{1}
\setlist[otazkyP]{label=\textbf{\sffamily\color{fisZelenaTmava}P\arabic*.}, ref=P\arabic*, resume,
  leftmargin=3em, itemsep=6pt, topsep=3pt}

% Výpisy z R (vstupují ze souborů vygenerovaných skriptem, nikdy ne opsané ručně)
\fvset{fontsize=\small, xleftmargin=0.5em, frame=leftline, rulecolor=\color{fisLinka},
  framerule=1.5pt, framesep=5pt}
% výpis se nedělí mezi stránky
\newcommand{\vystupR}[1]{\par\noindent\begin{minipage}{\linewidth}\VerbatimInput{#1}\end{minipage}\par}
